\renewcommand{\thechapter}{А}
\setcounter{chapter}{1}
\setcounter{section}{0}
\clearpage
\begin{flushright}
\textbf{ПРИЛОЖЕНИЕ А}
\end{flushright}

\begin{center}
\textbf{Структура кода и воспроизводимость}
\end{center}

\vspace{0.8cm}
\addcontentsline{toc}{chapter}{Приложение А. Структура кода и воспроизводимость}

\section{Структура репозитория}

Репозиторий разделён на несколько функциональных частей. В каталоге \texttt{src/} находятся функции построения геометрий, формирования признаков и базовых моделей. Каталог \texttt{tests/} содержит автоматические проверки геометрических процедур, генерации данных и вспомогательных функций моделей. В \texttt{data/} размещены сохранённые расчётные массивы, используемые в главах с численными результатами. Каталог \texttt{notebooks/} содержит воспроизводимые расчётные блокноты для физических проверок, сравнения моделей и анализа остатков. В \texttt{reports/} сохранены итоговые таблицы и графики, а каталог \texttt{thesis/} содержит исходные файлы настоящего текста.

\section{Воспроизводимость}

Основная проверка программной части выполняется командой \texttt{conda run -n diplom-kwant python -m pytest tests -q}.

Расчётные блокноты воспроизводятся командой \texttt{jupyter nbconvert --to notebook --execute notebook.ipynb --inplace} для соответствующего файла из каталога \texttt{notebooks/}. Такой порядок отделяет прямой расчёт и анализ данных от текста дипломной работы: численные таблицы и рисунки сначала формируются в \texttt{reports/}, затем используются в \LaTeX{}-документе.

\section{Окружение}

Расчёты выполнялись в окружении \texttt{diplom-kwant}. Для воспроизведения требуются Python-библиотеки для численных расчётов, обработки таблиц, построения графиков, машинного обучения и работы с Kwant. В частности, численная линейная алгебра и стандартные процедуры машинного обучения опираются на общепринятые реализации SciPy и scikit-learn~\cite{virtanen2020scipy,pedregosa2011sklearn}. Конкретный набор зависимостей фиксируется проектными файлами и текущим окружением; перед повторным запуском расчётов необходимо активировать то же окружение или создать совместимое.

\section{Тесты}

Автоматические тесты проверяют корректность геометрических построений, согласованность формирования расчётных наборов и поведение функций базовых моделей. Отдельные проверки относятся к физически мотивированным признакам, конфигурации гребневой регрессии и многослойного перцептрона, а также к валидации входных параметров. Эти тесты не заменяют физические проверки из основной части работы, но уменьшают риск технических ошибок в расчётной процедуре.

\section{Сгенерированные расчётные блокноты и отчёты}

Итоговые численные материалы сформированы в блокнотах этапов 06c, 07, 08 и 09. Этап 07 создаёт таблицу физических проверок и рисунки масштабирования, круговой проверки, вырождения \(dE_2\) и дискретизационных диагностик. Этап 08 формирует таблицы сравнения гребневой регрессии и многослойного перцептрона, включая посеменную устойчивость и значения по группам исключения. Этап 09 сохраняет точечные остатки гребневой регрессии, корреляции с диагностическими величинами и сводку проверки гипотезы об остатках. Соответствующие CSV-таблицы и PNG-рисунки находятся в \texttt{reports/} и \texttt{reports/assets/}.

\clearpage
